\section{Restatement of Purpose}

To whom it may concern.

Someone suggested there needs to be some ``dialog", presumably between Gornahoor and some unidentified dialogist. That
presumes there is something to discuss, necessitating some common ground, and resulting in a common goal. The ground is
Tradition. The goal is the recovery of Tradition in the West. We have not deviated from that in over five years.

In January, 2008, we set the stage in one of the first posts here on the Prophecies of \textbf{Rene Guenon}, Part 1\footnote{\url{https://gornahoor.net/p=35}} and
Part 2\footnote{\url{https://gornahoor.net/?p=36}}. Some commenters seem to believe we are trying to be persuasive although we have been clear that we are only
interested in facts, their consequences, and the various possibilities in give situation. These are the two relevant
facts.

\begin{itemize}
\item The West had a valid Tradition up until the Medieval Era 
\item The Tradition was eventually lost, after the destruction of the Templars according to Guenon, and perhaps a little
later in Evola's view 
\end{itemize}
The following corollary immediately follows from those two observations:

\begin{itemize}
\item In order to fully understand Tradition, it is sufficient to look to medieval Europe, its spiritual tradition, and
its social structure 
\end{itemize}
In other words, it is not necessary to look to alien traditions, but rather to Europe's
intellectual and spiritual past. This was brought to our attention by an essay by \textbf{Ananda Coomaraswamy}\footnote{\url{http://www.gornahoor.net/?p=556}}, in
which he made clear that a European could not understand the Vedanta unless he was able to understand the great
traditional writers of the West. There is a corollary to that observation:

\begin{itemize}
\item Anyone who scoffs at the spiritual tradition of medieval Europe is unlikely to recognize or understand the
traditional elements in the Vedanta. 
\end{itemize}
A project immediately presents itself, one which Gornahoor has endeavored to fulfill:

\begin{itemize}
\item Investigate the medieval spiritual tradition and demonstrate its teachings on metaphysics, psychology,
anthropology, history, sociology, politics, and so on. 
\end{itemize}
Now, all this is more than a mental exercise. The ultimate purpose is to point out a possible way for the recovery of
Tradition in the West. Moreover, that purpose is not the result of a merely arbitrary and subjective preference. As the
modern world cannot continue to sustain itself on its present path, Tradition will necessarily return. However, that
return can be relatively peaceful or quite painful. As noted in the links to Guenon's predictions,
there are three possibilities:

\begin{itemize}
\item Degeneration. If no effort is made to restore Tradition, the West will degenerate into barbarity. \textbf{Carl
Jung} noticed that in the psychology of his patients. Concerning the abandonment of the Western spiritual tradition by
modern men, he had this to say:

\begin{quotex}
At a time when a large part of mankind is beginning to discard Christianity, it may be worth our while to try to
understand why it was accepted in the first place. It was accepted as a means of escape from the brutality and
unconsciousness of the ancient world. As soon as we discard it, the old brutality returns in force … This is not a step
forwards, but a long step backwards into the past.
\end{quotex}
\item Assimilation. Outsiders from the East will bring Tradition back to Europe, whether by consent or by force. In the
former case, Europeans would adopt the Traditions of the newcomers. In the latter, the Tradition would come by force
marked by painful ethnic or racial conflicts. 
\item Transformation. In this scenario, the West returns to Tradition spontaneously and voluntarily. This would be
initiated by a group of men that Guenon calls the ``elite", in the sense that they understand tradition and ``true
intellectuality". They would be the initial leavening that would eventually spread. 
\end{itemize}
Since we are not interested in promoting the cause of degeneration, and since we are not ``outsiders",
Gornahoor's real project is almost too obvious to mention.

\begin{itemize}
\item To promote the spontaneous and voluntary transformation of the West. 
\end{itemize}
There are several preparatory steps required. Here are some we are engaged in:

\begin{itemize}
\item Make clear what the elite must know, both intellectually and morally. That is why we have provided so many
translations into English of rare works and have brought to light many ancient texts. Obviously, we read them in the
light of Tradition and not as moderns who have little conception of their real meaning. 
\item Encourage readers to engage with writings that are known to be initiatory, such as Dante, The Romance of the Rose,
Boccaccio, etc. 
\item Investigate spiritual and intellectual currents that still retain real and recoverable traces of tradition. 
\end{itemize}
There are other initiatives in motion, but it is too early to discuss them.

There is absolutely nothing new or novel in all this, it all derives from Guenon and other writers on Tradition. Hence,
there is really nothing to dialog about; you either accept the premise or you don't. We are
interested only in the former. It is not dialogists that are required, but rather those who want to develop such ideas
further, based on their own abilities and interests.



\flrightit{Posted on 2013-10-14 by Cologero }

\begin{center}* * *\end{center}

\begin{footnotesize}\begin{sffamily}



\texttt{jc00001 on 2013-10-14 at 09:37 said: }

The reintroduction of Neo-Platonism into Western thought would lay the foundations for a true living Tradition if
understood correctly. Plotinus's Enneads is the great masterwork of the Western Tradition.


\hfill

\texttt{Thorsten on 2013-10-14 at 10:34 said: }

Cologero,

thanks for the clarification. When I suggested a dialogue I was referring primarily to the fact that self-styled
``traditionalists" in the New Right/Alt Right need to wiling to be more open to traditional doctrine, especially in the
shape it assumed in Medieval Europe, and to be more open to metaphysical reality in general. (Otherwise they really
have no business calling themselves traditionalists.)

On the other hand I would like to humbly suggest that as prideful and sensually/aesthetically oriented [What else is
left in the absence of a living tradition than tribal and racial identity and old symbolic language (e.g. art) from a
bygone civilization? These seem to be the default surrogates for tradition, and the only remaining alternatives to
chaos when you were raised outside of tradition or have been heretofore fed such a degraded version of tradition that
you cannot in good conscience submit to it.] these apostates are that they are not all beyond saving as a group, and
should be handled as misguided rather than willfully contrarian. Now doubtless a number of the most strident voices in
these circles are in fact willfully contrarian, puerile, and deliberately disrespectful of Western tradition to a
degree that is shameful and they are probably not worth the time of Gornahoor or any other advocates of western
tradition.

I'm confident that most of the New Right/Alt Right are not interested in Muslim overlordship and
therefore are opposed the the option of Assimilation, although I'm not so sure that they would all
be opposed to some level brutality and barbarity in terms of a more militarized society and harsher corporal punishment
(It's undeniable that western society was markedly more brutal and barbarous until quite recently
in this regard.) If brutality and barbarism refers to coarseness of taste, vulgarity, and surrender to hedonism I think
its obvious that most of them would oppose Degeneration in this sense, as this is frequently the subject of articles of
the New Right. On the other hand I think the third possibility of Transformation might actually be received quite well
as this would appeal to those who were not born into an existing tradition or have been alienated the contemporary
Chrisitianity of any sort because they were raised on some anti-traditional low church protestantism which is devoid
traditional values such as inequality, metaphysics as opposed to literalist readings of scripture, and the primacy of
the intellect as opposed to blind faith.

I myself am not a member in any sense of the New Right, but I do think that there is more potential among those who
frequent such circles than there appears to be if you only focus on the stance of the writers and editors. As a former
Nietzsche-maniac, and stubborn materialist like so many of those frequent and comment on such websites I think that
many of these iconoclasts could be esotericists in the making.

Kind regards,

Thorsten


\hfill

\texttt{JA on 2013-10-14 at 10:54 said: }

when did Evola see the Tradition as being lost in Europe ? I was unaware he differed in chronology from Guenon.


\hfill

\texttt{cosmodromium on 2013-10-14 at 10:59 said: }

Gornahoor maintains an abnormally high quality of thought on clearly defined subject matter. Seeking ``dialogue" –
especially online – especially with the ``new/alt right" crowd would be diverting time and energy to lower the bar. I
believe somewhere Guenon states a viable mission for the man of tradition to bear witness to the modern age – as a
``watcher" – and maintain purity of doctrine. Which in my opinion Gornahoor has done admirably. Put the material out for
those with eyes to see. The others will invariably dilute, distort, or ``debate it. In fact, if I were to be so
obnoxious as to advance unsolicited advice, I would even suggest cutting loose the ballast of the ``comments" section.


\hfill

\texttt{jc00001 on 2013-10-14 at 11:25 said: }

Gornahoor is definitely the spiritual core of the counter-revolution and the living Tradition, online. The only site
that comes anywhere near is Ken Wheeler's \url{http://kathodos.com} and that sadly is private,
invitation only now.


\hfill

\texttt{Thorsten on 2013-10-14 at 12:02 said: }

cosmodromium,

Perhaps you're right about the new/alt right. Maybe I give them more credit than they deserve. My
opinion of them has already sunk quite a bit but I think that fostering an intellectual elite will be insufficient in
the long run for the project of Transformation. Unless I'm missing something I
can't help but conclude that some action must be taken at some point to reach the Kshatriya, i.e.
the politically inclined, in order to make some effect on the wider society. Perhaps I'm ``jumping
the gun" and need to be more patient but time is not on our side I'm afraid as the prospects of
Degeneration and Assimilation are materializing quite rapidly. (And I hope I don't come across as
arrogant. It's not my intention.) Again when I speak of dialogue I do not mean to suggest debating
doctrine or metaphysics, certainly not with those who are less informed but rather to brainstorm the logistics of how a
traditional society could be made to work at the present level of technological complexity, globalization,
``multiculturalism", mass media etc. It's a tedious and daunting project but I
don't see how it can be avoided.


\hfill

\texttt{Thomas Blanchard on 2013-10-14 at 14:01 said: }

Thorsten,

I'll chime in since I am loosely involved in various New Right spheres, and can echo most of your
feelings on this, although I think I've come to a different understanding of things.

I agree that the reformation of an authentic new second estate (warrior caste) will be a necessary step in an overall
return to order in the west, just as it will be necessary to reform the third estate on healthy principles.
Gornahoor's primary mission (and someone correct me if I'm wrong) is the
recovery of the western tradition and and the spiritual development of individuals who will eventually form a truly
legitimate first estate that will serve as its elite vanguard.

This being said, a parallel effort to build a true second estate does make sense, but they key is that it should be
properly subordinated to a new first estate; the castes and vocations should have proper hierarchical relations with
one another. All of the various New Right movements of today represent different manifestations of this impulse, but
without a proper spiritual authority to guide them (and you are correct that the modern exoteric Church as the world
sees it is rather weak and uninspiring), these efforts constitute a rebellious attempt to re-establish the Bronze Age
(written of by Hesiod). In this case I respectfully disagree with Cologero's non-distinction
between pure modernists and the New Right types that one might find at a place like Counter-Currents: modernists
represent the Iron Age, the ``Renaissance", the French Revolution, and the ``Enlightenment", but the New Right truly
wishes for the return of the Bronze Age and their writers are men like Homer and Nietzsche.

In a sense, the dialogue you have been talking about has been taking place for a long time – you can see it in
Plato's works: Socrates' debates with Thucydides or especially with
Callicles, where his opponents represent these more ``New Right" sorts.

As a person who also has mixed proclivities and tends to sympathize with Nietzsche, I share your feelings, but
it's best for the time being to focus on re-establishing a spiritual elite before trying to give
direction to the chaotic mess of political agitators who are united only in their opposition to the Iron Age.

For your personal development, if you cannot commit yourself to a religion now, at least I would recommend taking up
things like traditional martial arts and mountain climbing – Evola wrote primarily for people of our type, and some of
his last advice to young people was not to worry about forming an ``order", but to simply focus on building character. I
believe that this is good advice.


\hfill

\texttt{Thorsten on 2013-10-14 at 20:06 said: }

Thomas,

I'm on the same page with you and Gornahoor on the primacy of establishing a reformed
spiritual/intellectual elite/oratores/brahmin caste (call it what you will). Thought precedes action and the wise must
rule the strong and I say that is all the more reason to quash counterfeit spiritualities propagated by some in the new
right and to formulate a working alternative. How to do this specifically I am not sure and won't
presume to tell anyone.

I myself cannot say with certainty whether I am more a bellator than an orator but I do have some experience with rock
climbing and I intend to take up some form of martial arts when I have the means to do so. Anyway I understand the
necessity of rebuilding from the top down as has been reiterated here at Gornahoor many times, but if we
don't get things rebuilt (bring about the Transformation) before society collapses into utter
degeneracy do we simply allow some semblance of order to be imposed on the West by Wahhabi radicals (Assimilation) by
means of brutality and barbarism which may make the barbarian incursions in the last days of the Roman empire look
merciful in comparison?

To be clear am no enemy of Islam. I am only pointing out that when the present ideology of materialism and
egalitarianism implodes upon itself the most consistent (and forceful) ideology in the West will not only be Muslim but
the most ferocious variant thereof.


\hfill

\texttt{Michael on 2013-10-14 at 20:25 said: }

Cologero, I understand the concept of an elite, but I don't understand how the elite will act as a
``leaven" to transform society. Is it purely through contemplation? Or will there also be action similar to the way that
enlightenment ideas spread through society? From reading Guenon and Tomberg, I have the impression that it is an unseen
spiritual leavening process.


\hfill

\texttt{Cologero on 2013-10-14 at 23:18 said: }

@ Thomas Blanchard Those are mostly excellent suggestions … more generally, besides rock climbing, there are other
elements to master depending on your personal equation. See \textit{Qualifications for Initiation}\footnote{\url{http://www.gornahoor.net/?p=1161}} for example. If only the NR
would follow your advice and stay out of religious and spiritual matters. Their proper domain is the manifestation
which should be the focus. An important point you bring up is the ``third estate", which does not get the attention it
needs. The modern world is in the hands of that estate, not some decadent spiritual tradition as the NR seems to be
obsessed with.

There is some subtleties with Socrates, who was executed for atheism and impiety. His call for the rule of philosophers
in touch with the ideas replaced the Homeric vision of the rule of the gods. So you could see him as the precursor of
modernity on the one hand, or else as someone who spread esoteric teachings prematurely.


\hfill

\texttt{Cologero on 2013-10-14 at 23:32 said: }

@Thorsten

I put up the quotes from Alice Bailey to show how events happen on their own accord apart from the wills of men. She saw
turmoil inaugurating a new age of one religion, one race, etc, although it is hard to be so sanguine about that. It is
odd that the Theosophists seen eternal progress instead of the cycles of decline, considering their reliance on Asian
teachings.

Back to the point, events will unfold as they do. Decisions made and actions executed today will have consequences down
the road. I don't know how Guenon was able to see the possibly imminent ethnic conflict in 1921,
but they are certainly foreseeable now. Yet it doesn't have to be that way.

I don't know why you feel any personal agitation or anxiety. It is not as though you are in charge
of anything nor is the worst likely to happen in your lifetime. Besides rock climbing, seek out
Seneca's 41st letter to Lucillius:

\begin{quotex}
If you see a man who is unterrified in the midst of dangers, untouched by desires, happy in adversity, peaceful amid the
storm, who looks down upon men from a higher plane and view the gods on a footing of equality, will not a feeling of
reverence for him steal over you? …

When a soul rises superior to other souls, when it is under control, when it passes through every experience as if it
were of small account, when it smiles at our fears and at our prayers, it is stirred by a force from heaven. A thing
like this cannot stand upright unless it be propped by the divine. 

\end{quotex}
That is the ideal to strive for. Jung's commentary on that passage is also quite interesting, but
I'll save it for another time.


\hfill

\texttt{Cologero on 2013-10-14 at 23:44 said: }

Guenon ties it to a very specific event: the destruction of the Templars. The reason for this may be difficult to
understand or even to accept, but it has to do both with the exodus of esoteric orders from Europe and the ties between
esoteric Christian and esoteric Islamic orders. Guenon has left a very interesting clue about that, although I am not
ready to make it public without understanding it more deeply. I have pointed out the clue on Gornahoor although it will
take some digging.

I don't believe Evola dates it so precisely since he did not have the same interest in that topic
as Guenon. Instead he refers to the regression of the castes over a longer time. As we have often pointed out, Evola
said he believed what every well-bred person prior to the French Revolution considered normal and healthy.


\hfill

\texttt{Cologero on 2013-10-14 at 23:47 said: }

Yes, jc, and neoplatonic authors are well represented on Coomaraswamy's list. Unfortunately,
professional philosophers today consider it an obsolete phase of their discipline. Theologians are worse when the say
something like ``early theologians used Greek philosophy as the vehicle to express their theology", as though it were a
totally arbitrary choice and any modern philosophy would serve the same purpose.


\hfill

\texttt{Cologero on 2013-10-14 at 23:54 said: }

I don't understand either, Michael, as the Spirit will blow where he wills. Yes, false ideas spread
quite easily, but true ideas are longer lasting. There is always the feeling of having to ``do something" but instead we
need that interior silence, not unlike what Seneca described, which connects us to Heaven which allows the influx of
forces.


\hfill

\texttt{Logres on 2013-10-15 at 06:06 said: }

``gain when I speak of dialogue I do not mean to suggest debating doctrine or metaphysics, certainly not with those who
are less informed but rather to brainstorm the logistics of how a traditional society could be made to work at the
present level of technological complexity, globalization, ``multiculturalism", mass media etc. It's
a tedious and daunting project but I don't see how it can be avoided."

Thorsten, this may be no reason not to spend time on it now, but the ``present level of complexity" may decrease quite
drastically over the next 20-50 years. At least on so large a scale.


\hfill

\texttt{Thorsten on 2013-10-15 at 11:52 said: }

@ Cologero,

I'm fairly well acquainted with the Stoics, including Seneca, but I'm not
asking for advice on how to be more equanimous as my emotions are actually quite subdued to the point of being a bit
atrophied largely by virtue of studying Stoicism and practicing meditation. If anything I've found
it too easy to become detached from the world to the point that I see no point in taking any action. But such a level
of detachment seems to be almost antithetical to the Indo-European tradition. So in the interest of taking action on
earth beyond my own head I am compelled to insist on answers to a number of riddles about the application of
traditional principles on that which is happening right here right now. Therefore the urgency which I show is not
indicative of anxiety or agitation (A warrior may be calm and fearless in battle but he still must be swift and
decisive and for this he needs a battle plan and tactics.) and at least b) and c) below are hardly more my predicaments
than they are all of ours whether we choose to confront them or not. If they do not affect us personally they will
affect our descendants. Rather the urgency is indicative of the fact that I (and I would have to assume others as well)
remain without any unambiguous answers to and am therefore unable to plan and make decisions with regards to the
problems of:

a) What is my tradition if I was not born into one? And if we are to work toward Transformation, at what point do we
abandon our current exoteric religions (assuming one is already attached to a religion)? Do we remain faithful by all
external appearances to our religion, all the while building a new vehicle for tradition and then as it were suddenly
apostasize from the religion that we currently follow when the new vehicle is ready? (whether this be in 5 years or
200) I understand Evola's stance about Riding the Tiger and not being compelled to get involved in
politics or religion, but apostasy from religion of any sort is inconsistent with the advice of Guenon et al. is it
not?

b) In order to prepare the groundwork for the Transformation do we simply, learn to ``watch our thoughts", ``stop
thinking", cultivate ``interior silence" and other such contemplative exercises then build a community of like-minded
brahmanic individuals, rather than ``do something" out in the real world in terms of metapolitical cultural activities
(influencing the media, education etc.) and political activism? Are such activities the business of the second estate?
If the Transformation is to be spontaneous how passive (in the conventional sense of the word) must we be in terms of
creating the new institutions. Should we have any sense of urgency whatsoever? Or is it so important to ensure that the
new vehicle of tradition is fully formed that matters of politics and demographic change should be of no concern no
matter how tumultuous things may become if we do not act? I ask because I think it is indeed possible to make a
difference if enough individuals care to.

c.) On that note I think it is worth pointing out that men of tradition, at least those of the second estate, have
always fought until the bitter end against foreign invaders no matter how traditional the religion of their foes may
have been and no matter what the odds were. Just how is our situation any different? (Byrhtwold encapsulates the
Indo-European attitude in such a situation. The same attitude that a stoic Roman, Spartan, or Vedic warrior showed.
\url{http://en.wikiquote.org/wiki/The\_Battle\_of\_Maldon}) If anything the battle will be largely a political rather than
the nasty kind if enough pressure can be brought to bear on governments in the next decade. Is the degeneracy of the
contemporary West reason enough to surrender to a more traditional invading enemy? Is the value of ethnic loyalty
trumped by the interest of submitting to a more traditional social order? (It bears mentioning radical Muslims are
generally opposed to the authority of the first estate).

Again I pose these questions not as rhetorical devices or as incitement but out of a sincere interest in knowledge of
the application of Tradition to our present circumstances. And if you are not comfortable providing answers I will
understand.

P.S.

@ Logres,

I think such a simplification would be the best possible scenario. 

Thanks,

Thorsten


\hfill

\texttt{Avery on 2013-10-15 at 13:41 said: }

A frightening possibility is that the world will not be simplifying anytime soon, but will continue to confuse one or
two generations more with astounding quantities and technological wonders, at the price of any last scraps of
tradition. In that scenario both degeneration and assimilation become highly likely. Yet Gornahoor makes transformation
sound so easy. Considering the dangers involved it is actually the most difficult of the three paths.

Seraphim Rose probably did more for this transformative process than any of us. He was completely orthodox and he seems
to have had prophetic abilities as well. Yet one of his followers started a cult in his name outside the communion of
the ROCOR (in Brookline, MA). Can we really say that what we revive will be the West?


\hfill

\texttt{Cologero on 2013-10-15 at 20:36 said: }

Thorsten, you have no standing to ``insist" on anything here. I have been providing answers in several hundred thousand
words of original material and translations without feeling the least discomfort. If you want to address specific
points discussed, then give it a try, otherwise it is time for you to move on.


\hfill

\texttt{Thorsten on 2013-10-15 at 21:03 said: }

Cologero,

I meant no disrespect. I suppose the word ``insist" was not the right term. Anyway I thought I had addressed specific
points which you have discussed but in the interest of not creating any bad blood I will cease to ask these questions.

Best Regards,

Thorsten


\hfill

\texttt{Logres on 2013-10-16 at 06:33 said: }

Avery, that's a good observation – the next generation or two will be crucial, \& the likely real
test, because they do have enough resources to string it along for a little further. However, at least to my thinking,
it's getting more and more obvious that the group that is benefiting is shrinking and shrinking,
and this is in our favor, if we can strengthen the classes that are being shoved aside to keep the profit margins thick
enough. This is why a new ``elite" is paramount. The situation is conducive for opening eyes to the disgusting nature of
materialism.


\hfill


\end{sffamily}\end{footnotesize}
